\section{Conclusion}\label{sec:conclusion}
We compared four rules for authorizing academic credentials within one smart-contract design. At a common 4-of-7 threshold, requiring a registrar, a dean and an examination officer did not increase the number of insiders needed. It excluded \BlockCoalExcluded{} of the 35 coalitions of that size and cost \GasRoleOverhead{} additional gas per issuance. It also reduced the number of absences that can block issuance from four to two, and lowered simulated completion from \OutageFiveC{} to \OutageFiveD{} under a 5\% whole-role outage. Additional holders of each role reduced the loss from individual absence and not from whole-role outages. The benefit of roles was modest under independent compromise and large when collusion was correlated within offices. A single issuer is cheap but fails with one insider. Unanimity resists every coalition smaller than the whole group but is rarely available.

Role constraints are therefore best seen as a way to require collusion across offices and not as a substitute for a higher threshold, and they should be accompanied by enough holders of each role. Future work should remove reliance on one governor, study conditional delegation, replace assumed collusion patterns by evidence, replace assumed delays by observed ones, require several parties for revocation, and test the system on a public or permissioned network. None of these changes would by itself establish that the source record is true.
